\documentclass[10pt,a4paper]{amsart}
\usepackage[margin=19mm]{geometry}
\usepackage{amsmath,amssymb,mathtools}
\usepackage[scr=rsfs,cal=euler]{mathalfa}
\usepackage{xfrac}
\usepackage[unicode,hidelinks,bookmarks=false]{hyperref}
\newcommand{\ie}{i.e.\ }
\newcommand{\cf}{cf.\ }
\newcommand{\resp}{respectively\ }
\newcommand{\Subsection}{\subsection}
\providecommand{\nobreakdash}{\nobreak\hbox{-}}
\setlength{\emergencystretch}{2em}
\multlinegap0pt
\usepackage{bm}
\usepackage{bbm}
\usepackage{dsfont}
\usepackage{xspace}
\usepackage{cancel}
\usepackage{mathtools}
\usepackage{amsthm}
\makeatletter
\@ifundefined{theorem}{\newtheorem{theorem}{Theorem}}{}
\@ifundefined{lemma}{\newtheorem{lemma}[theorem]{Lemma}}{}
\@ifundefined{notation}{\newtheorem{notation}[theorem]{Notation}}{}
\makeatother
\title[Normal 2-coverings in affine groups of small dimension]{Normal 2-coverings in affine groups of small dimension}
\author{Marco Fusari, Andrea Previtali, Pablo Spiga}
\date{Source: arXiv:2507.15610v1 (CC BY 4.0)}
\begin{document}
\maketitle

\noindent\emph{Source and licence.} Normal 2-coverings in affine groups of small dimension. Version arXiv:2507.15610v1 (CC BY 4.0), distributed under the Creative Commons Attribution 4.0 International licence, as recorded in the arXiv licence metadata for this exact version and shown on the abstract page https://arxiv.org/abs/2507.15610v1. Full rights notice: licenses/arxiv-2507.15610-CC-BY-4.0.txt.\par\smallskip

\noindent\textit{Editorial scope.} Counted unit: the Theorem whose source label is \texttt{thrm:mainmain} in the article (source0.tex lines 383--386, source label \texttt{thrm:mainmain}), delivered together with the article's own complete original proof of it (source0.tex lines 387--427, one \texttt{proof} environment of 41 lines). No mathematics is written, repaired or re-derived: statement and proof are the article's own text line for line. Required context retained uncounted: lemma:mercoledi0 (lines 195--197); not (lines 207--230); lemma:mercoledi2 (lines 249--253); lemma:mercoledi4 (lines 335--343); lemma:navarro (lines 370--372). Every definition, display and label of those lines is retained unchanged. Omitted from this package and not redistributed: 1--194, 198--206, 231--248, 254--334, 344--369, 373--382, 428--755. Nothing inside a retained excerpt is omitted or altered: every retained body is delivered exactly as the article's lines are written, so no omission receipt of any kind accompanies this package. Citations: the retained excerpts carry no citation command at all, so no bibliography entry of the article is transcribed and no reference is redistributed. Reference adaptation: none was needed and none was made; every reference inside the delivered unit resolves inside it, so no unresolved reference or citation is produced. Printed slips kept literal: no printed word, symbol, hypothesis or number of the counted statement or of its proof was altered. Layout and class: the article's own class is \texttt{amsart} with packages including xcolor, amsmath, amsfonts, amssymb, color, hyperref, enumerate, tabularx, amsthm, amscd, verbatim, graphicx, multirow, tikz; the wrapper is the lane's pinned \texttt{amsart} editorial wrapper at 10pt on A4, which restates the theorem-like environments the retained text needs and adds the article's own macro definitions that the retained text needs (none), with extra wrapper packages (bm, bbm, dsfont, xspace, cancel, mathtools). The delivered numbering is the unit's own. Assets: no retained span contains a figure environment, a graphics-inclusion command, a PDF-page inclusion, a table or a TikZ picture; no image file is copied into this package, the package declares no asset dependency and no image file is redistributed with it. Licence: version arXiv:2507.15610v1 of this article is distributed under the Creative Commons Attribution 4.0 International licence, as recorded in the arXiv licence metadata for this exact version and shown on the abstract page https://arxiv.org/abs/2507.15610v1. A full rights notice is delivered at licenses/arxiv-2507.15610-CC-BY-4.0.txt. Characters: every retained excerpt is pure ASCII, so the delivered engine, XeLaTeX with its default Latin Modern text font, reports no missing character.
\begin{lemma}\label{lemma:mercoledi0}
Let $G=V\rtimes H$ be an affine primitive group with $H$ nilpotent. Then $G$ is non-basic (that is, $\gamma(G)\ne 2$) if and only if $H=\langle H^\ast\rangle$.
\end{lemma}

\begin{notation}\label{not}{\rm
Let $p$ be a prime number, let $f$ be a positive integer, let $q = p^f$, let $\mathbb{F}_q$ be the finite field of cardinality $q$, let $\lambda$ be a generator of the multiplicative group of $\mathbb{F}_q$, and let $\sigma$ be the Frobenius automorphism of $\mathbb{F}_q$. Thus $x^\sigma = x^p$ for all $x \in \mathbb{F}_p$.

With this notation, we have:
\begin{align*}
\Gamma\mathrm{L}_1(q) &= \mathbb{F}_q^\ast \rtimes \mathrm{Aut}(\mathbb{F}_q) = \langle \lambda \rangle \rtimes \langle \sigma \rangle,\\
\mathrm{A}\Gamma\mathrm{L}_1(q) &= \mathbb{F}_q \rtimes \Gamma\mathrm{L}_1(q) = \mathbb{F}_q \rtimes \langle \lambda, \sigma \rangle.
\end{align*}
Note that $\Gamma\mathrm{L}_1(q)$ is metacyclic, and hence so is each of its subgroups.

Let $H$ be a subgroup of $\Gamma\mathrm{L}_1(q)$. In particular, we may write $H = \langle \lambda^d, \sigma^j \lambda^i \rangle$, where $d$ divides $q - 1$, $j$ divides $f$, and $i \in \mathbb{Z}$. We may assume that $\langle \lambda^d \rangle = H \cap \langle \lambda \rangle$. Since $|H : \langle \lambda^d \rangle| = |\langle \sigma^j \rangle|$, it follows that
$$(\sigma^j\lambda^i)^{f/j} = \lambda^{i\frac{p^f - 1}{p^j - 1}} \in H \cap \langle \lambda \rangle = \langle \lambda^d \rangle,$$
that is,
\begin{equation}\label{eq:neq1}
d\, \text{ divides }\, i\frac{p^f - 1}{p^j - 1}.
\end{equation}

An arbitrary element $h$ of $H$ can be written (uniquely) as $h=(\sigma^j\lambda^i)^k\lambda^{d\ell}$, for a unique $k\in \{0,\ldots,f/j-1\}$ and a unique $\ell\in \{0,\ldots,(p^f-1)/d-1\}$. We have
\begin{equation}\label{often}
h=(\sigma^j\lambda^i)^k\lambda^{d\ell}=\sigma^{ijk}\lambda^{i\frac{p^{jk}-1}{p^j-1}+d\ell}.
\end{equation}
We use very often this computation in this section.
}
\end{notation}

\begin{lemma}\label{lemma:mercoledi2}
Let $H = \langle \sigma^j \lambda^i, \lambda^d \rangle$ be as in Notation~$\ref{not}$. If $H = \langle H^\ast \rangle$, then $p^j - 1$ divides both $d$ and $i$.

In particular, replacing $H$ by a suitable $\langle \lambda \rangle$-conjugate, we may assume $H = \langle \sigma^j,\lambda^d \rangle$.
\end{lemma}

\begin{lemma}\label{lemma:mercoledi4}
Let $H$ be as in Notation~$\ref{not}$ and assume $H=\langle H^\ast\rangle$ and $H$  nilpotent. 

If $p=2$, or if $p^j\equiv 1\pmod 4$, or if $f/j$ is odd, then $(p^f-1)/d$ divides $f/j$.
Moreover, when $p^j\equiv 3\pmod 4$ and $f/j$ is even, then $(p^f-1)/d$ divides
$(f/j)2^{v_2(p^j+1)-v_2(d)}$.

In particular, in all cases $((p^f-1)/d)_{2'}$ divides $f/j$.
\end{lemma}

\begin{lemma}\label{lemma:navarro}
Let $J$ be a normal subgroup of a finite group $H$ with $H/J$ abelian. Let $V$ be an irreducible $\mathbb{F}_pH$-module and let $W$ be an irreducible $\mathbb{F}_pJ$-submodule of $V$. Then $\dim V/\dim W$ is an integer dividing $|H:J|$.
\end{lemma}

\begin{theorem}\label{thrm:mainmain}
Let $G=V\rtimes H$ be a primitive subgroup of $\mathrm{A}\Gamma\mathrm{L}_1(q)$
 with $H$ nilpotent and let $r$ be the largest prime dividing $|H|$. If $r>2$, then  $r$ divides $|H:\langle H^\ast\rangle|$; in particular, $\gamma(G)=2$ and $G$ is basic.
\end{theorem}

\begin{proof}
Let $K=\langle H^\ast\rangle$.
As $H^\ast\subseteq K$, we have $K^\ast=H^\ast$ and hence $K=\langle K^\ast\rangle$. Using Notation~\ref{not}, we have
$H=\langle\sigma^j\lambda^i,\lambda^d\rangle$. Moreover, applying Notation~\ref{not} and Lemma~\ref{lemma:mercoledi2} to $K$, replacing $K$ and $H$ with a suitable $\langle\lambda\rangle$-conjugate, we may suppose $K=\langle \sigma^{j_0},\lambda^{d_0}\rangle$.

We argue by contradiction and we suppose $|H:K|$ is not divisible by $r$, where $r>2$ is the largest prime divisor of $|H|$. As $|H|=(f/j)(p^f-1)/d$ and $|K|=(f/j_0)(p^f-1)/d_0$, we deduce that
\begin{equation}\label{eq:today2}
|H:K|=\frac{j_0}{j}\cdot \frac{d_0}{d}\,\hbox{ is relatively prime to }r.
\end{equation}

Assume  $r\in \pi(f/j)\setminus\pi((p^f-1)/d)$. Let $R$ be a Sylow $r$-subgroup of $H$ and let $Q$ be a Hall $r'$-subgroup of $H$. As $H$ is nilpotent, $H=R\times Q$. Since $|H:K|$ is relatively prime to $r$, we have $R\le K$ and, as $\gcd(r,(p^f-1)/d_0)=1$, we have $R\le \langle \sigma^{j_0}\rangle$ and $R$ is abelian. In particular,
$$H\le {\bf C}_{\Gamma\mathrm{L}_1(p^f)}(R)\le {\bf C}_{\Gamma\mathrm{L}_1(p^f)}(\sigma^{f/r})=\mathbb{F}_{p^{f/r}}^\ast\rtimes \langle \sigma\rangle.$$
Observe now that the subfield $\mathbb{F}_{p^{f/r}}$ is $H$-invariant, because it is invariant by $\mathbb{F}_{p^{f/r}}^\ast\rtimes \langle \sigma\rangle$. However, this contradicts the fact that $H$ acts irreducibly on $\mathbb{F}_{p^f}$. This contradiction has shown that $r\in\pi((p^f-1)/d)$ and hence, from~\eqref{eq:today2}, we get
\begin{equation}\label{eq:allprimes}
r\in \pi((p^f-1)/d_0)\cap \pi(f/j_0).
\end{equation}

 From~\eqref{eq:allprimes}, $r\in \pi((p^f-1)/d_0)$ and hence, from Lemma~\ref{lemma:mercoledi4} applied to $K$, we have
\begin{align}\label{eq:today3}
v_r(f/j_0)\ge v_r((p^f-1)/d_0).
\end{align}


Now consider  $J=\langle \sigma^{j_0},\lambda^d\rangle$ and  $W=\mathbb{F}_p[\lambda^d]$, the subfield of $\mathbb{F}_{p^f}$ containing $\lambda^d$. Clearly, $K\le J\le H$ and $J\unlhd H$. Observe that $W$ is an $\mathbb{F}_pJ$-module because $W$ is invariant under $\lambda^d$ and $\sigma^{j_0}$. Moreover, $W$ is irreducible as $\mathbb{F}_pJ$-module from its definition. Since $W$ is a subfield of $\mathbb{F}_{p^f}$, we have  $W=\mathbb{F}_{p^{f_0}}$ for some divisor $f_0$ of $f$. Observe that $f_0$ is the smallest positive integer with ${\bf o}(\lambda^d)=(p^f-1)/d$ dividing $p^{f_0}-1$. In other words, $f_0={\bf o}_{(p^f-1)/d}(p)$ and, by raising $p$ to the $j_0$th power, we get $f_0/\gcd(f_0,j_0)={\bf o}_{(p^f-1)/d}(p^{j_0})$. From Euler's theorem,
\begin{equation}\label{eq:today}\frac{f_0}{\gcd(f_0,j_0)} \hbox{ divides }\varphi((p^f-1)/d).
\end{equation}

Observe now that $J\unlhd H$ and $|H:J|$ is relatively prime to  $r$. Therefore, from Lemma~\ref{lemma:navarro}, $f/f_0$ is relatively prime to $r$. Thus, from~\eqref{eq:today2} and~\eqref{eq:today}, 
\begin{equation}\label{eq:today1}
v_r(f/j)\le v_r(\varphi((p^f-1)/d)).
\end{equation}
From~\eqref{eq:today2},~\eqref{eq:today3} and~\eqref{eq:today1}, we have
\begin{align}\label{eq:ultimate}
v_r(f/j_0)&\ge v_r((p^f-1)/d_0)=v_r((p^f-1)/d)>v_r(\varphi((p^f-1)/d))\ge v_r(f/j),
\end{align}
where the strict inequality 
$v_r((p^f-1)/d)>v_r(\varphi((p^f-1)/d))$
follows from the fact that $r$ is the largest odd prime dividing $|H|$. However, the chain of inequalities in~\eqref{eq:ultimate} gives a contradiction.

Finally, as $H\ne \langle H^\ast\rangle$, from Lemma~\ref{lemma:mercoledi0}, $\gamma(G)=2$ and $G$ is basic.
\end{proof}

\end{document}
